%! @@TITLE@@ — Unit @@UNITINT@@, Lesson @@LESSONID@@ — Lesson Plan
% Gradual-release plan (I do / we do / you do together). Teacher-facing — use the formal
% vocabulary and the lettered standard codes freely. There is NO spoiler rule in this course.
% This course HAS `fixedskillbox` (unlike AP Statistics): use it where a tabularx must stay
% intact on one page, and plain `skillbox` + \boxguard everywhere else.
\documentclass[10pt]{article}
\usepackage{@@PREFIX@@-article}
\usepackage{@@PREFIX@@-boxes}
\usepackage{graphicx}   % -article does NOT load graphicx; needed for thumbnails/screenshots
\graphicspath{{images/}}
\newcommand{\TallMath}[1]{$\displaystyle #1\rule[-1.4em]{0pt}{3.2em}$}
@@COURSEMACROS@@\newcommand{\UnitNumberName}{Unit @@UNITINT@@: @@UNITTITLE@@ \quad}
\newcommand{\LessonNumberName}{Lesson @@LESSONID@@: @@TITLE@@}

\begin{document}
\begin{center}
    {\Large\bfseries \CourseName \\
    {\small\bfseries \UnitNumberName \LessonNumberName}}
\end{center}

\begin{tcolorbox}[colback=frost, colframe=cerulean, boxrule=0.9pt, arc=2mm,
    left=2mm, right=2mm, top=1.5mm, bottom=1.5mm]
    \textbf{Primary Objective:} TODO --- what students can do, interpret, and \emph{justify}
    by the end.
    \\[2pt]
    % CITE THE STANDARD CODES. This is a standards-driven course; the codes appear here and
    % again on the Reinforcement & Extension / Connections line.
    \textbf{Standards:} TODO (e.g.\ PS.DS.1a, PS.DS.1c)
\end{tcolorbox}

\renewcommand{\arraystretch}{1.4}
\begin{skillbox}[Priority Ideas \& Skills]{goldbox}
% Fill in the \item text ONLY. Two tabularx cells, one itemize each — do not re-nest or
% re-balance the environments (that is the classic source of a stray \end{itemize}).
\begin{tabularx}{\linewidth}{@{}X|X@{}}
\textbf{Learning targets (student language)} & \textbf{Key understandings (the ``why'')} \\[2pt]
\hline
\vspace{-6pt}
\begin{itemize}[leftmargin=1.1em, nosep, after=\vspace{2pt}]
  \item TODO: ``I can \dots'' target, traceable to a lettered Knowledge \& Skill (cite the code).
\end{itemize}
&
\vspace{-6pt}
\begin{itemize}[leftmargin=1.1em, nosep, after=\vspace{2pt}]
  \item TODO: the why --- and state the lesson's \textbf{target misconception} explicitly.
\end{itemize}
\end{tabularx}
\end{skillbox}

\begin{skillbox}[{Vocabulary, Concepts, \& Theorems}]{greenbox}
    % TODO: key terms + definitions. These are TAUGHT, in guided-notes order, and they appear
    % on the cover, in the notes vocabbox, and on the slides. Nothing is withheld.
\end{skillbox}

% fixedskillbox never splits, so the phase table stays intact on one page.
% THE PHASE TABLE IS A CONTRACT: the Lesson box's per-section minutes must sum to 20, and the
% activity must be workable in 18. If a phase does not fit, cut it — do not let this table lie.
\begin{fixedskillbox}[Lesson at a Glance --- \MeetingLength]{frost}
\renewcommand{\arraystretch}{1.25}
\begin{tabularx}{\linewidth}{@{}l c X X@{}}
\textbf{Phase} & \textbf{Min} & \textbf{Students} & \textbf{Teacher} \\
\hline
Warm-Up & 5 & Three spiral items, silent, then check with a neighbor & Circulate; debrief one item aloud \\
Guided Notes \& Practice & 20 & Hook, then fill notes sections 1--TODO with the class; work the practice box & \emph{I do / we do} --- model section TODO, run the rest as a class, release the practice box \\
Group Activity & 18 & TODO activity title, items 1--TODO, groups of 3 & \emph{You do together} --- circulate, question rather than answer, choose work to display \\
Debrief & 7 & Defend answers; correct their own in a second color & Run the displayed work; land TODO; whole-class check for understanding \\
Close \& Assign & 5 & Record the assignment; preview next lesson & Name what changed today; assign packet \emph{or} DeltaMath \\
\end{tabularx}
\end{fixedskillbox}

\begin{fixedskillbox}[Activate Prior Knowledge \& Spiral Review]{frost}
    \begin{tabularx}{\linewidth}{@{}X c@{}}
        \begin{minipage}[t]{0.55\linewidth}\vspace{0pt}
            % TODO: the ~3 warm-up items, and for each one WHICH NOTES SECTION picks it up
            % again — the warm-up should read as a running start, not a detour.
        \end{minipage}
        &
        \begin{minipage}[t]{0.38\linewidth}\vspace{0pt}\centering
@@SPIRALWARMUP@@
        \end{minipage}
    \end{tabularx}
\end{fixedskillbox}

\boxguard
\begin{skillbox}[Hook]{frost}
    TODO: what goes on the board before anyone sits down, and the question to ask. When the
    lesson's crux is a misconception, state the claim, take three or four answers, write them
    down without judging any of them, and \textbf{leave it unresolved} --- name which notes
    section settles it.
\end{skillbox}

\boxguard
\begin{skillbox}[Lesson --- Guided Notes (20 min)]{frost}
\begin{multicols}{2}
% ONE bolded paragraph per notes section, EACH CARRYING ITS OWN MINUTE BUDGET. The budgets
% must sum to 20 including the practice box. Say what to model (I do), what to fill in with
% the class (we do), and what to release (you do).
\textbf{1. TODO (N min).} TODO.

\textbf{2. TODO (N min). \emph{This is the I-do section.}} TODO --- model it explicitly.

\textbf{3. TODO (N min). \emph{We do.}} TODO. If this section sets up the crux, take a vote,
record the split, and \textbf{do not settle it here}.

\textbf{4. TODO (N min). \emph{The lesson.}} TODO --- the crux. It lands only because the class
already committed to a side.

\textbf{Practice box (N min, \emph{you do}).} Release it. TODO: what the arithmetic already
gives them, so the time goes to the classification and the interpretation sentence.
\end{multicols}
\end{skillbox}

\boxguard
\begin{skillbox}[Explicit Instruction: TODO the core move]{frost}
\begin{tabularx}{\linewidth}{@{} p{0.46\linewidth} X @{}}
\begin{minipage}[t]{\linewidth}\vspace{0pt}
\textbf{Steps --- give these verbatim}
\begin{enumerate}[leftmargin=1.4em, itemsep=2pt, topsep=2pt]
  \item TODO.
\end{enumerate}
\end{minipage}
&
\begin{minipage}[t]{\linewidth}\vspace{0pt}
\textbf{Worked example --- TODO}
\begin{itemize}[leftmargin=1.1em, itemsep=2pt, topsep=2pt]
  \item[1.] TODO.
\end{itemize}
\textit{Contrast: TODO --- the near-miss case that looks the same and is not.}
\end{minipage}
\end{tabularx}
\end{skillbox}

\boxguard
\begin{skillbox}[Active Monitoring]{redbox}
\begin{itemize}[leftmargin=1.2em, nosep, itemsep=2pt]
  \item \textbf{Watch for TODO.} Keyed to a notes section AND an activity item number.
  \item \textbf{Watch for reasons that are not reasons.} TODO: the two phrasings to reject
        all period, and what to require instead.
\end{itemize}
Cold-call prompts: TODO.
\end{skillbox}

\boxguard
\begin{skillbox}[Group Work (18 min)]{redbox}
\small
\textbf{One activity, the whole class, groups of three.} Everyone works items 1--TODO off the
same data set; the gold \emph{Extension} box at the end is for groups that get there, not a
separate assignment. Hand it out and start the clock --- there is nothing to assign or sort.

\vspace{2pt}
\textbf{The arc.} TODO: which items are the fast on-ramp (readable straight off the data),
which compute and demand a sentence naming the group and the units, \textbf{which item is the
crux}, and which is the hardest thing on the page.

\vspace{2pt}
\textbf{Circulate with questions, not answers.} \emph{Item N:} ``TODO prompt.''
\emph{Item N, if they are stuck at TODO:} ``TODO prompt.''

\vspace{2pt}
\textbf{Pick the work to display for the debrief:} TODO --- name the specific items, and get
two \emph{disagreeing} answers for the crux.
\end{skillbox}

\boxguard
\begin{skillbox}[Debrief (7 min) --- what goes on the board]{frost}
Run the displayed group work in this order. \textbf{Students correct their own answers in a
second color} as you go; that self-correction is the record you collect from.
\begin{enumerate}[leftmargin=1.3em, nosep, itemsep=3pt]
  \item \textbf{Item N.} TODO.
  \item \textbf{Item N --- the crux.} TODO: put both displayed answers up and land the rule
        out loud, in one sentence.
  \item \textbf{Extension item N}, if any group reached it --- otherwise pose it aloud. TODO.
\end{enumerate}

\vspace{2pt}
\textbf{Check for understanding --- ask it cold, whole class.} TODO the question. Hands down,
thirty seconds of think time, then cold-call three students. Correct answer: TODO.
\emph{This replaces the exit ticket.}

\vspace{2pt}
\textbf{Formative read.} Sort the answers into three piles: \textbf{(a)} TODO correct and
justified; \textbf{(b)} TODO correct for the wrong reason --- the big, fixable middle;
\textbf{(c)} TODO missed it. Open the next lesson from pile (b).
\end{skillbox}

\boxguard
\begin{skillbox}[Reinforcement \& Extension]{goldbox}
\textbf{Homework --- TODO context.} $\sim$6 items in new contexts spanning this lesson's lettered
skills and spiraling one earlier idea, each ending in an interpretation or a justification.
TODO: what each item targets, how it is scored, and when it is due. \textbf{Scored} --- the
cover's score column carries a blank for it.

\textbf{Packet or DeltaMath --- decide this lesson.} The packet set is authored either way.
DeltaMath's statistics coverage is thin, so the packet is the default; assign the DeltaMath set
instead only where it genuinely carries TODO (this lesson's skills). Say which one aloud at
Close \& Assign so it is unambiguous.

\textbf{Extension (optional).} TODO.

\textbf{Preview of the next lesson.} TODO.

\textbf{Connections \& Big Ideas.} \textit{Standards covered:} TODO. \textit{Spirals forward
to:} TODO. \textit{Reaches back to:} TODO.
\end{skillbox}

% ── Teacher notes — one per component, in packet order ──
% TEACHERNOTES: teacher prose lives HERE, never in a _key. A note in a key is the one block
% with no counterpart in the blank, so it makes the key run longer and costs the student
% packet a blank page. FOUR notes — there is none for the debrief, which is fully specified
% in its own box above.
\begin{teachernote}[Warm-Up]
    % TODO: pacing; which item separates the class; what to cut if it runs long.
\end{teachernote}

\begin{teachernote}[Guided Notes]
    % TODO: the per-section pacing again, in prose; which section IS the lesson and which to
    % compress if you are behind; the single highest-leverage move (usually: set up the crux
    % and refuse to settle it early); what the practice box's pre-worked arithmetic buys you.
\end{teachernote}

\begin{teachernote}[Group Activity]
    % TODO: groups of three, 18 minutes, one activity for everyone — nothing to assign.
    % The 10-minute mark and where groups should be; which item every group must have
    % attempted; what to do if the room runs slow; what to do if nobody reaches the Extension.
\end{teachernote}

\begin{teachernote}[Homework]
    % TODO: what to flag when assigning (the items that cannot be answered by computing);
    % which item to grade first; whether DeltaMath covers this lesson; which item to go over
    % at the start of the next lesson.
\end{teachernote}
\end{document}
